\documentclass[12pt,a4paper]{article}
\usepackage[utf8]{inputenc}
\usepackage[T1]{fontenc}
\usepackage{amsmath,amssymb,amsfonts}
\usepackage{geometry}
\usepackage{booktabs}
\usepackage{hyperref}
\geometry{margin=2.5cm}

\title{Differential Predictions of the Causal Horizon Framework}
\author{Kevin Mu\~noz}
\date{}

\begin{document}
\maketitle

\begin{abstract}
Steps 1--4 of the Causal Horizon Framework (CHF) derived the primordial scalar power spectrum entirely from the geometry and thermodynamics of the apparent horizon, closing the four analytical gaps of the formalization document. The resulting tilt formula is $n_s - 1 = -2\varepsilon + 2\beta_S/3$, with $\varepsilon = \frac{1}{2}d\ln S/d\ln a$ (horizon entropy rate) and $\beta_S = 0$ for the minimal conformally coupled field on the two-dimensional horizon. This document identifies the observables by which the CHF differs from standard single-field inflation and specifies the conditions under which those differences are measurable. Four differential predictions are established. First, in the minimal case ($\beta_S = 0$), the CHF gives $n_s - 1 = -2\varepsilon$ --- missing the second slow-roll contribution $\eta'$ that standard inflation allows as an independent degree of freedom. For the same background $\varepsilon$, the two frameworks predict tilts differing by $\eta'$; future CMB experiments with $\sigma(n_s) \sim 0.002$ will be sensitive to this at the level of $|\eta'| \sim 0.01$. Second, the standard consistency relation $r = 16\varepsilon$ was confirmed in Step~6 by deriving tensor modes from the horizon boundary, so the CHF predicts a specific line in the $(n_s, r)$ plane: $r = 8(1 - n_s)$, corresponding to $r \approx 0.28$ for the Planck 2018 central value $n_s = 0.9649$. This is in tension with the current bound $r < 0.056$. With $r = 16\varepsilon$ confirmed by Step~6, the tension is genuine: the minimal framework ($\beta_S = 0$) is excluded by Planck 2018, and $\beta_S < -0.042$ is required. Third, the running $\alpha_s = -d^2\ln S/dN^2$ contains no $\xi^2$ term from a cubic potential. Fourth, the minimal CHF action is quadratic in the boundary field, predicting $f_{NL}^{\rm CHF} = 0$ at leading order from the horizon. These predictions are derived, not assumed, and provide the empirical handles by which the CHF can be discriminated from --- or falsified against --- the standard inflationary paradigm.
\end{abstract}

\section{Introduction}

A theoretical framework earns empirical weight by predicting observables differently from existing frameworks. Steps 1--4 of the CHF reproduced results already obtained in standard inflation (scale-invariant spectrum, tilt formula, boundary--bulk relation) from a different conceptual starting point --- the horizon, not the bulk inflaton. Matching known results is necessary for consistency but not sufficient for discrimination.

This document turns to the question of falsifiability: what would the CHF predict differently from a bulk inflationary model operating on the same background? Section~2 establishes the parameter structure of the CHF and compares it to that of standard inflation. Sections 3--6 derive the four differential predictions. Section~7 summarizes the current observational status. Section~8 records how Step~6 resolved Prediction~1.

\section{Parameter Structure: CHF vs.\ Standard Inflation}

\subsection{Standard single-field slow-roll inflation}

In standard single-field slow-roll inflation, the scalar power spectrum is determined by two independent slow-roll parameters at leading order:
\begin{equation}
n_s - 1 = -2\varepsilon - \eta', \qquad r = 16\varepsilon \tag{1}
\end{equation}
where $\varepsilon = -\dot{H}/H^2$ and $\eta'$ is a second slow-roll parameter encoding the rate of change of $\varepsilon$ and/or the curvature of the inflaton potential. The two parameters are independent: for fixed $\varepsilon$, $\eta'$ can range over a wide interval without violating any consistency condition. In the Hubble slow-roll hierarchy, $\eta' = \dot{\varepsilon}/(H\varepsilon) = d\ln\varepsilon/dN$.

This two-parameter freedom means that, for a given observed tilt $n_s$, the slow-roll parameter $\varepsilon$ --- and hence the tensor-to-scalar ratio $r = 16\varepsilon$ --- is not determined without additional information about $\eta'$.

\subsection{The CHF scalar sector}

The CHF (Steps 1--3) gives:
\begin{equation}
n_s - 1 = -2\varepsilon + \frac{2\beta_S}{3}, \qquad \varepsilon = \frac{1}{2}\frac{d\ln S}{d\ln a} \tag{2}
\end{equation}
where $\beta_S = m_\text{eff}^2/H^2$ is the normalized effective mass of the horizon field. For the minimal horizon field ($\beta_S = 0$, the unique choice with no intrinsic coupling to the two-sphere curvature):
\begin{equation}
n_s - 1 = -2\varepsilon \qquad (\text{minimal CHF}) \tag{3}
\end{equation}
This is a single-parameter relation: given the observed $n_s$, $\varepsilon$ is uniquely fixed to $(1-n_s)/2$. The second slow-roll parameter $\eta'$ that appears in standard inflation is absent from the CHF at this order.

The structural origin of this difference: in standard inflation, the mode variable is $u_k = z\mathcal{R}_k$ with $z = a\sqrt{2\varepsilon}\,M_{\rm Pl}$, and the effective potential $z''/z$ involves both $\varepsilon$ and $\eta'$ at leading slow-roll order. In the CHF, the mode variable is $u_k = z_A\psi_k$ with $z_A = a/\sqrt{\varepsilon}$, and the effective potential $z_A''/z_A$ involves $\varepsilon$ at leading order but the $\eta'$ correction enters at a different subleading level. For power-law inflation ($\eta' = 0$), the two mode equations are equivalent. For general slow-roll ($\eta' \neq 0$), the CHF and the bulk description predict different tilts for the same background.

\section{Prediction 1: The $(n_s, r)$ Plane}

\subsection{The CHF locus (confirmed by Step~6)}

Step~6 derived the tensor power spectrum from the horizon boundary field and confirmed the standard consistency relation $r = 16\varepsilon$ (see Step~6, Section~8). Equation (3) then implies:
\begin{equation}
r_{\rm CHF}(\beta_S = 0) = 16\varepsilon = 8(1 - n_s) \tag{4}
\end{equation}
This is a line in the $(n_s, r)$ plane, parameterized by $\varepsilon$ alone. For the Planck 2018 central value $n_s = 0.9649$:
\begin{equation}
r_{\rm CHF} = 8 \times 0.0351 = 0.281 \tag{5}
\end{equation}
Standard inflation, by contrast, spans a two-dimensional region in the $(n_s, r)$ plane. The CHF ($\beta_S=0$) collapses this to a single point.

\subsection{Tension with current data}

The Planck 2018 bound is $r < 0.056$ (95\% CL). The CHF prediction $r \approx 0.281$ lies well above this bound. Step~6 confirmed $r = 16\varepsilon$, so the tension is genuine and has two possible resolutions:

\textbf{(a) The minimal CHF ($\beta_S = 0$) is falsified.} With $r = 16\varepsilon$ confirmed, $\beta_S = 0$ is observationally excluded at 95\% CL by the combination $(n_s, r)_{\rm Planck}$.

\textbf{(b) The horizon field is non-minimal ($\beta_S \neq 0$).} For $\beta_S \neq 0$, the tilt is $n_s - 1 = -2\varepsilon + 2\beta_S/3$. With $r = 16\varepsilon$, for fixed $n_s$:
\begin{equation}
r = 8(1 - n_s) + \frac{16\beta_S}{3} \tag{6}
\end{equation}
Requiring $r < 0.056$ for $n_s = 0.9649$ gives:
\begin{equation}
\frac{16\beta_S}{3} < 0.056 - 0.281 = -0.225 \qquad \Longrightarrow \qquad \beta_S < -0.042 \tag{7}
\end{equation}
A slightly tachyonic horizon field ($m_\text{eff}^2 < 0$) is required. A tachyonic mass reduces the effective $\nu$ and contributes a blue shift to the tilt, allowing a given $n_s$ to be achieved with smaller $\varepsilon$ (and hence smaller $r$).

\subsection{Discrimination from the standard locus}

In the $(n_s, r)$ plane, the CHF prediction (4) lies on the line $r = 8(1-n_s)$. Standard slow-roll inflationary models populate a broad region below this line (since $\eta' < 0$ for models with concave potentials). The CHF line passes through or above most standard model predictions (Starobinsky, natural inflation, $\alpha$-attractors) that satisfy the Planck 2018 bound. A future measurement of $r$ above the current bound (e.g., $r \sim 0.01$--$0.05$) would be consistent with the $\beta_S < 0$ CHF but not with the minimal $\beta_S = 0$ case.

\section{Prediction 2: Absence of the Second Slow-Roll Parameter}

\subsection{The $\eta'$ correction}

In standard inflation, the tilt receives a contribution from $\eta' = d\ln\varepsilon/dN$:
\begin{equation}
\left(n_s - 1\right)_{\rm std} = -2\varepsilon - \eta' \tag{8}
\end{equation}
In the CHF ($\beta_S = 0$):
\begin{equation}
\left(n_s - 1\right)_{\rm CHF} = -2\varepsilon \tag{9}
\end{equation}
For the same cosmological background (same $\varepsilon(t)$), the difference is:
\begin{equation}
\Delta(n_s - 1) \equiv \left(n_s - 1\right)_{\rm CHF} - \left(n_s - 1\right)_{\rm std} = \eta' = \frac{d\ln\varepsilon}{dN} \tag{10}
\end{equation}
This difference vanishes for power-law inflation ($\varepsilon = \text{const}$, $\eta' = 0$) but is nonzero for all models where the Hubble rate is not decelerating at a constant rate.

\subsection{Magnitude and observational reach}

For typical inflationary models, $|\eta'| \sim |\varepsilon| \sim 0.01$--$0.03$. The Planck 2018 measurement has $\sigma(n_s) \approx 0.004$; future experiments (Simons Observatory, CMB-S4, LiteBIRD) are expected to reach $\sigma(n_s) \sim 0.002$. The difference $|\Delta(n_s-1)| = |\eta'| \sim 0.01$ is potentially detectable at the $\sim 5\sigma$ level with next-generation CMB data. If $r$ is measured (fixing $\varepsilon = r/16$), then $\eta'$ is the residual $\eta' = -(n_s - 1) - 2\varepsilon$ in standard inflation, while the CHF predicts this residual to vanish.

\subsection{Physical content}

In the CHF, the tilt has a single thermodynamic origin: $n_s - 1 = -d\ln S/d\ln a$. The rate of entropy change of the horizon encodes the full tilt. No second parameter is needed because the CHF treats $\varepsilon$ as the fundamental degree of freedom, fixing its relation to the spectrum through the geometry of the mode equation for $z_A = a/\sqrt\varepsilon$. Standard inflation allows an additional deformation through the curvature of the inflaton potential, producing the $\eta'$ term. The CHF prediction $\eta' = 0$ is equivalent to the statement that \textbf{the inflaton potential is exponential} (power-law inflation), which is the unique case where the CHF and standard inflation agree exactly.

\section{Prediction 3: Spectral Running from Horizon Thermodynamics}

\subsection{CHF running formula}

The running of the spectral index is defined as $\alpha_s \equiv dn_s/d\ln k \approx d(n_s-1)/dN$. From equation (3):
\begin{equation}
\alpha_s = \frac{d(n_s-1)}{dN} = -2\frac{d\varepsilon}{dN} = -2\varepsilon\,\eta_\varepsilon \tag{11}
\end{equation}
where $\eta_\varepsilon \equiv d\ln\varepsilon/dN$. In thermodynamic language, using $\varepsilon = \frac{1}{2}d\ln S/dN$:
\begin{equation}
\alpha_s = -\frac{d^2\ln S}{dN^2} \tag{12}
\end{equation}
The running of the spectral index equals \textbf{minus the second derivative of the horizon entropy} with respect to the number of e-folds.

\subsection{Relation to the standard running}

In standard inflation at leading order: $\alpha_s = -2\varepsilon\eta_H - \xi^2$, where $\xi^2 = M_{\rm Pl}^4 V'V'''/(V^2)$ involves the third derivative of the potential. The CHF running (11) matches the first term but has no analogue of $\xi^2$, since the horizon field action has no cubic potential. This predicts:
\begin{equation}
\alpha_s^{\rm CHF} = -2\varepsilon\eta_\varepsilon = (1-n_s)\,\eta_\varepsilon \tag{13}
\end{equation}
For $n_s = 0.965$ and $\eta_\varepsilon \sim -0.01$: $\alpha_s \approx 3.5\times10^{-4}$, consistent with Planck 2018 ($\alpha_s = -0.0045 \pm 0.0067$). A measurement of $|\alpha_s|$ significantly larger than $(1-n_s)|\eta_\varepsilon|$ would be inconsistent with the minimal CHF.

\subsection{Consistency relation}

Combining (9) and (13) for the minimal CHF:
\begin{equation}
\frac{\alpha_s}{(n_s - 1)^2} = \frac{-2\varepsilon\eta_\varepsilon}{4\varepsilon^2} = -\frac{\eta_\varepsilon}{2\varepsilon} = -\frac{d\ln\varepsilon/dN}{d\ln S/dN} \tag{14}
\end{equation}
This ratio is a pure horizon quantity. If the entropy grows at a constant rate ($\eta_\varepsilon = 0$, power-law inflation), both $\alpha_s$ and this ratio vanish exactly.

\section{Prediction 4: Vanishing Boundary Non-Gaussianity}

\subsection{The horizon field action is quadratic}

The minimal CHF action for the horizon field $\phi(\Omega, t)$ is (eq.~(23) of Document~5):
\begin{equation}
S_\phi = \int dt\, d\Omega\, R_A^2 \left[(\partial_t\phi)^2 + H^2(\nabla_\Omega\phi)^2\right] \tag{15}
\end{equation}
This action is \textbf{purely quadratic} in $\phi$: there are no cubic or higher self-interaction terms. The quantization of a quadratic action produces a Gaussian vacuum state. Therefore, the three-point function $\langle\psi_k\psi_{k'}\psi_{k''}\rangle$ vanishes at leading order from the boundary field alone:
\begin{equation}
f_{NL}^{\rm CHF,\, boundary} = 0 \tag{16}
\end{equation}

\subsection{Standard inflation prediction}

In standard single-field slow-roll inflation, non-Gaussianity is generated by self-interactions of the inflaton in the bulk. The leading result is (Maldacena 2003):
\begin{equation}
f_{NL}^{\rm std} = \frac{5}{12}(n_s - 1) + \frac{5}{6}\varepsilon + \frac{5}{12}\eta' \approx \mathcal{O}(\varepsilon, \eta') \sim -0.02 \tag{17}
\end{equation}
This is small but nonzero. The CHF predicts a contribution from the boundary that is identically zero.

\subsection{Observational status and discriminability}

The predicted difference $|f_{NL}^{\rm CHF} - f_{NL}^{\rm std}| \sim 0.02$ is far below the current Planck 2018 bound $|f_{NL}| < 5$ and below the sensitivity of forthcoming experiments (CMB-S4 aims for $\sigma(f_{NL}) \sim 1$). In the foreseeable observational future, both frameworks predict effectively Gaussian primordial fluctuations. A detection of $f_{NL} \sim \mathcal{O}(1)$ would be inconsistent with both.

\section{Observational Status}

\begin{table}[h]
\centering
\small
\begin{tabular}{@{}lp{2.8cm}p{2.8cm}p{2.6cm}p{3cm}@{}}
\toprule
Prediction & CHF ($\beta_S=0$) & Standard inflation & Current constraint & Status \\
\midrule
$r$ given $n_s = 0.9649$ & $r = 0.281$ & $r \in [0, 0.28]$ (2D family) & $r < 0.056$ & \textbf{Tension confirmed} --- $r = 16\varepsilon$ from Step~6; requires $\beta_S < -0.042$ \\[4pt]
$\eta'$ contribution to tilt & $\eta' = 0$ & $\eta'$ free & $\sigma(n_s) = 0.004$ & Not yet discriminated \\[4pt]
$\alpha_s$ & $-2\varepsilon\eta_\varepsilon$, no $\xi^2$ & $-2\varepsilon\eta_H - \xi^2$ & $\alpha_s = -0.0045 \pm 0.0067$ & Consistent \\[4pt]
$f_{NL}$ (boundary) & $0$ & $\mathcal{O}(0.01)$ & $|f_{NL}| < 5$ & Consistent, not discriminable \\
\bottomrule
\end{tabular}
\end{table}

The only prediction currently in tension is the value of $r$ from the minimal CHF. Step~6 confirmed $r = 16\varepsilon$, establishing that the minimal framework ($\beta_S = 0$) is excluded by Planck 2018. Consistency with current data requires $\beta_S < -0.042$.

\section{Resolution by Step~6 (Tensor Perturbations)}

Prediction~1 (the $(n_s, r)$ locus) depended on whether the CHF satisfies the standard consistency relation $r = 16\varepsilon$. Step~6 derived the tensor power spectrum from the two TT polarization modes on the apparent horizon boundary.

\textbf{Key result of Step~6:} The tensor pump field is $z_T = a$ (vs.\ $z_A = a/\sqrt{\varepsilon}$ for scalars). Tensor modes are perturbations of the \textit{shape} of the horizon (traceless transverse shear), carrying no slow-roll suppression, while scalar modes perturb the \textit{size} ($\psi = \delta R_A/R_A$), suppressed by $\sqrt\varepsilon$. The calculation yields:
\begin{equation}
\Delta_h^2 = \frac{2H^2}{\pi^2 M_{\rm Pl}^2}, \qquad n_T = -2\varepsilon, \qquad r = 16\varepsilon \tag{18}
\end{equation}
The standard consistency relations $r = -8n_T$ and $r = 16\varepsilon$ are confirmed.

\textbf{Consequence for Prediction~1:} With $r = 16\varepsilon$ confirmed, the minimal CHF ($\beta_S = 0$) predicts $r \approx 0.281$ for $n_s = 0.9649$, above the Planck 2018 bound $r < 0.056$. The minimal framework is excluded at 95\% CL.

\textbf{CHF-specific relation:} Step~6 also established an additional consistency relation specific to the CHF (for $\beta_S = \beta_T = 0$):
\begin{equation}
n_T = n_s - 1 \qquad \Longrightarrow \qquad r = 8(1 - n_s) \tag{19}
\end{equation}
This is a differential prediction of the CHF relative to standard inflation, where $n_T$ and $n_s - 1$ are independent at the same slow-roll order.

\section{Summary}

\begin{table}[h]
\centering
\small
\begin{tabular}{@{}lp{3.5cm}p{2.5cm}p{4cm}@{}}
\toprule
Prediction & Expression & Observable & Discriminating power \\
\midrule
$(n_s, r)$ locus & $r = 8(1-n_s)$, confirmed $r=16\varepsilon$ (Step~6) & CMB $B$-modes & High --- minimal CHF excluded at 95\% CL by Planck 2018 \\[4pt]
Absence of $\eta'$ & $(n_s - 1)_{\rm CHF} - (n_s-1)_{\rm std} = \eta'$ & Future $\sigma(n_s) \sim 0.002$ & Moderate --- accessible with CMB-S4/LiteBIRD \\[4pt]
Running from thermodynamics & $\alpha_s = -d^2\ln S/dN^2$, no $\xi^2$ & CMB multipole shape & Low --- $\xi^2$ contribution too small \\[4pt]
Boundary non-Gaussianity & $f_{NL}^{\rm boundary} = 0$ & CMB bispectrum & Very low --- $|f_{NL}| \sim 0.01$ below experimental reach \\
\bottomrule
\end{tabular}
\end{table}

The framework's most immediate empirical handle is the value of $r$. The minimal CHF ($\beta_S = 0$) predicts $r \approx 0.28$; Step~6 confirmed $r = 16\varepsilon$, establishing this prediction as genuine and the minimal framework as excluded at 95\% CL by Planck 2018. Consistency with current data requires $\beta_S < -0.042$. The next discriminating observable is the absence of $\eta'$ in the tilt, accessible with future CMB experiments at $\sigma(n_s) \sim 0.002$.

\begin{thebibliography}{9}
\bibitem{planck2018} Planck Collaboration, ``Planck 2018 results. X. Constraints on inflation,'' \textit{A\&A} \textbf{641}, A10 (2020).
\bibitem{maldacena2003} J. M. Maldacena, ``Non-Gaussian perturbations in inflationary cosmologies,'' \textit{JHEP} \textbf{0305}, 013 (2003).
\bibitem{baumann2009} D. Baumann, ``TASI Lectures on Inflation,'' arXiv:0907.5424 (2009).
\bibitem{munoz2025a} K. Mu\~noz, ``A Horizon-Centered Formalization of Cosmological Dynamics and Perturbations,'' Causal Horizon Framework, Document~5 (2025).
\bibitem{munoz2025b} K. Mu\~noz, ``Scale Invariance from Horizon Field Quantization,'' Causal Horizon Framework (2025).
\bibitem{munoz2025c} K. Mu\~noz, ``Gauge-Invariant Formulation of the Horizon Perturbation,'' Causal Horizon Framework (2025).
\bibitem{munoz2025d} K. Mu\~noz, ``Spectral Tilt from Horizon Thermodynamics,'' Causal Horizon Framework (2025).
\bibitem{munoz2025e} K. Mu\~noz, ``The Boundary--Bulk Relation from the Gauss--Codazzi Equations,'' Causal Horizon Framework (2025).
\bibitem{munoz2025f} K. Mu\~noz, ``Tensor Perturbations from the Horizon Boundary,'' Causal Horizon Framework (2025).
\end{thebibliography}

\end{document}
